% Paper 3 of the RH-Q3 Trilogy
% Title: Weil positivity and the Riemann Hypothesis via operator modules
% Author: Eugen Malamutmann
% Date: January 2026

\documentclass[11pt,a4paper]{article}

% Packages
\usepackage{amsmath,amssymb,amsthm}
\usepackage{mathtools}
\usepackage{enumitem}
\usepackage[margin=2.5cm]{geometry}
\usepackage{hyperref}
\usepackage[capitalize]{cleveref}

% Theorem environments
\newtheorem{theorem}{Theorem}[section]
\newtheorem{lemma}[theorem]{Lemma}
\newtheorem{proposition}[theorem]{Proposition}
\newtheorem{corollary}[theorem]{Corollary}
\theoremstyle{definition}
\newtheorem{definition}[theorem]{Definition}
\newtheorem{remark}[theorem]{Remark}

% Notation shortcuts
\newcommand{\RR}{\mathbb{R}}
\newcommand{\ZZ}{\mathbb{Z}}
\newcommand{\TT}{\mathbb{T}}
\newcommand{\CC}{\mathbb{C}}
\newcommand{\WW}{\mathcal{W}}
\newcommand{\GG}{\mathcal{G}}

% Title information
\title{Weil positivity and the Riemann Hypothesis\\ via operator modules}
\author{Eugen Malamutmann\thanks{University of Duisburg--Essen. ORCID: 0000-0003-4624-5890}}
\date{January 2026}

\begin{document}

\maketitle

\begin{abstract}
We complete the operator-theoretic proof of the Riemann Hypothesis via Weil's positivity criterion. Building on Part~I (Fej\'er--heat density and Lipschitz control) and Part~II (Toeplitz barrier and RKHS prime contraction), we establish that the Weil quadratic functional $Q$ is nonnegative on the full Weil cone $\WW$ of even, compactly supported test functions. The proof proceeds by combining the uniform A3 bridge inequality $\lambda_{\min}(T_M[P_A] - T_P) \ge c_*/4 > 0$ with the density of Fej\'er--heat generators and the Lipschitz continuity of $Q$: positivity on generators propagates to each compact window $\WW_K$, and the union over $K$ yields $Q \ge 0$ on $\WW$. Applying Weil's classical equivalence, this positivity implies that all nontrivial zeros of the Riemann zeta function lie on the critical line $\Re s = 1/2$. The entire argument is analytic: every bound is established from explicit inequalities, the parameters are given in closed form, and no numerical tables or K-dependent schedules enter the main proof chain.
\end{abstract}

%==============================================================================
\section{Introduction}\label{sec:intro}
%==============================================================================

This paper completes a trilogy developing an operator-theoretic proof of the Riemann Hypothesis (RH) via Weil's positivity criterion. The Weil criterion, established in~\cite{Weil1952,Weil1972}, provides a remarkable equivalence: the location of the nontrivial zeros of the Riemann zeta function---whether they all lie on the critical line or not---is encoded in the positivity of a single quadratic form.

\begin{theorem*}[Weil's Positivity Criterion]
The Riemann Hypothesis holds if and only if $Q(\Phi) \ge 0$ for all admissible test functions $\Phi$.
\end{theorem*}

Our approach verifies this positivity through a modular chain of analytic estimates:
\begin{center}
\begin{tabular}{c}
$\mathrm{(T0)} + \mathrm{(A1')} + \mathrm{(A2)} + \mathrm{(A3)} + \mathrm{(RKHS)}$ \\
$\Downarrow$ \\
$Q(\Phi) \ge 0$ for all $\Phi \in \WW$ \\
$\Downarrow$ \\
Riemann Hypothesis
\end{tabular}
\end{center}

The components of this chain are:
\begin{itemize}
\item \textbf{(T0) Normalization}: The Guinand--Weil crosswalk between frequency conventions, established in this paper.
\item \textbf{(A1$'$) Density}: The Fej\'er$\times$heat cone $\GG_K$ is dense in the cone of continuous, even, nonnegative functions on $[-K,K]$, established in Part~I~\cite{Paper1}.
\item \textbf{(A2) Continuity}: The functional $Q$ is Lipschitz continuous on each compact window, established in Part~I~\cite{Paper1}.
\item \textbf{(A3) Toeplitz barrier}: The uniform Archimedean floor $\min P_A(\theta) \ge c_* = 11/10$, established in Part~II~\cite{Paper2}.
\item \textbf{(RKHS) Prime cap}: The uniform bound $\|T_P\| \le \rho(t_{\mathrm{rkhs}}) < c_*/4$, established in Part~II~\cite{Paper2}.
\end{itemize}

\subsection*{Main result}

\begin{theorem}[Riemann Hypothesis]\label{thm:RH-intro}
Under the analytic module chain $\mathrm{(T0)} + \mathrm{(A1')} + \mathrm{(A2)} + \mathrm{(A3)} + \mathrm{(RKHS)}$, all nontrivial zeros of the Riemann zeta function $\zeta(s)$ lie on the critical line $\Re s = 1/2$.
\end{theorem}

The proof combines the uniform A3 bridge inequality from Part~II with a density-continuity argument: positivity on Fej\'er--heat generators implies positivity on each compact window $\WW_K$, and the union over $K$ gives positivity on the full Weil cone $\WW$. Weil's criterion then converts this functional positivity into the zero-location statement.

\subsection*{What is new}

Three features distinguish this approach:
\begin{enumerate}
\item \textbf{Uniform bounds}: All constants ($c_*$, $t_{\mathrm{sym}}$, $t_{\mathrm{rkhs}}$, $B_{\min}$) are fixed globally and independent of the compact window $[-K,K]$.
\item \textbf{No numerical tables}: Every bound is derived from explicit analytic estimates (digamma asymptotics, Gaussian integrals, Gershgorin bounds).
\item \textbf{Modular architecture}: Each component (density, continuity, Toeplitz barrier, RKHS cap) is self-contained and verifiable independently.
\end{enumerate}

\subsection*{Organization}

Section~\ref{sec:T0} establishes the Guinand--Weil normalization (T0). Section~\ref{sec:modules} summarizes the analytic module stack from Parts~I--II. Section~\ref{sec:closure} proves the main positivity theorem on $\WW$. Section~\ref{sec:weil} applies the Weil criterion to deduce the Riemann Hypothesis. Section~\ref{sec:discussion} discusses the broader context and verification philosophy.

%==============================================================================
\section{The Guinand--Weil Normalization}\label{sec:T0}
%==============================================================================

We work on the frequency axis with the standard Fourier transform convention
\begin{equation}\label{eq:fourier}
\widehat{\varphi}(\xi) = \int_{\RR} \varphi(t)\,e^{-2\pi i t\xi}\,dt, \qquad \varphi(t) = \int_{\RR} \widehat{\varphi}(\xi)\,e^{2\pi i t\xi}\,d\xi.
\end{equation}

\subsection{The Weil functional}

\begin{definition}[Archimedean density]\label{def:arch-density}
The Archimedean density is
\[
a(\xi) := \log\pi - \Re\psi\Bigl(\frac{1}{4} + i\pi\xi\Bigr),
\]
where $\psi = \Gamma'/\Gamma$ is the digamma function. The scaled density is $a^*(\xi) := 2\pi\,a(\xi)$.
\end{definition}

\begin{definition}[Prime nodes and weights]\label{def:prime-nodes}
The prime sampling nodes are
\[
\xi_n := \frac{\log n}{2\pi}, \qquad n \ge 2,
\]
and the Weil weights are
\[
w_Q(n) := \frac{2\Lambda(n)}{\sqrt{n}},
\]
where $\Lambda$ is the von Mangoldt function.
\end{definition}

\begin{definition}[Weil functional]\label{def:weil-functional}
For an even, compactly supported test function $\Phi$ with $\mathrm{supp}\,\Phi \subset [-K, K]$, the Weil quadratic functional is
\begin{equation}\label{eq:Q-def}
Q(\Phi) := \int_{\RR} a^*(\xi)\,\Phi(\xi)\,d\xi - \sum_{n \ge 2} w_Q(n)\,\Phi(\xi_n).
\end{equation}
\end{definition}

Since $\Phi$ has compact support, the integral reduces to $[-K,K]$ and the sum is finite over $2 \le n \le e^{2\pi K}$.

\subsection{Normalization crosswalk}

The following proposition establishes that our normalization matches the classical Guinand--Weil form.

\begin{proposition}[Guinand--Weil matching]\label{prop:GW-match}
Under the change of variables $\eta = 2\pi\xi$, the functional $Q$ defined by~\eqref{eq:Q-def} coincides with the classical Guinand--Weil functional
\[
Q_{\mathrm{GW}}(\varphi_{\mathrm{GW}}) := \int_{\RR} \Bigl(\log\pi - \Re\psi\Bigl(\frac{1}{4} + \frac{i\eta}{2}\Bigr)\Bigr)\,\varphi_{\mathrm{GW}}(\eta)\,d\eta - \sum_{n \ge 2} \frac{\Lambda(n)}{\sqrt{n}}\,\bigl(\varphi_{\mathrm{GW}}(\log n) + \varphi_{\mathrm{GW}}(-\log n)\bigr),
\]
where $\varphi_{\mathrm{GW}}(\eta) = \Phi(\eta/(2\pi))$.
\end{proposition}

\begin{proof}
The substitution $\eta = 2\pi\xi$ gives $d\eta = 2\pi\,d\xi$ and $\psi(1/4 + i\eta/2) = \psi(1/4 + i\pi\xi)$. For the Archimedean integral:
\[
\int_{\RR} \Bigl(\log\pi - \Re\psi\Bigl(\frac{1}{4} + \frac{i\eta}{2}\Bigr)\Bigr)\,\varphi_{\mathrm{GW}}(\eta)\,d\eta = \int_{\RR} 2\pi\,a(\xi)\,\Phi(\xi)\,d\xi = \int_{\RR} a^*(\xi)\,\Phi(\xi)\,d\xi.
\]
For the prime sum, $\varphi_{\mathrm{GW}}(\pm\log n) = \Phi(\pm\xi_n)$, and by evenness $\Phi(\xi_n) + \Phi(-\xi_n) = 2\Phi(\xi_n)$, yielding
\[
\sum_{n \ge 2} \frac{\Lambda(n)}{\sqrt{n}}\,\bigl(\varphi_{\mathrm{GW}}(\log n) + \varphi_{\mathrm{GW}}(-\log n)\bigr) = \sum_{n \ge 2} \frac{2\Lambda(n)}{\sqrt{n}}\,\Phi(\xi_n).
\]
Combining gives $Q(\Phi) = Q_{\mathrm{GW}}(\varphi_{\mathrm{GW}})$.
\end{proof}

\begin{remark}[Normalization invariance]
The positivity of $Q$ is independent of the choice of Fourier normalization: different conventions lead to rescaled densities that preserve the sign of $Q(\Phi)$.
\end{remark}

%==============================================================================
\section{The Analytic Module Stack}\label{sec:modules}
%==============================================================================

We summarize the analytic inputs established in Parts~I--II of this series.

\subsection{Module (A1$'$): Fej\'er--heat density}

\begin{definition}[Fej\'er--heat window]
For $B \ge 1$ and $t > 0$, the Fej\'er--heat window is
\[
\Phi_{B,t}(\xi) := \Bigl(1 - \frac{|\xi|}{B}\Bigr)_+ e^{-4\pi^2 t\,\xi^2}.
\]
The Fej\'er--heat cone $\GG_K$ consists of finite nonnegative linear combinations of even symmetrizations of such windows, restricted to $[-K,K]$.
\end{definition}

\begin{theorem}[Density, Part~I]\label{thm:density}
For every $K > 0$, the cone $\GG_K$ is dense in $C^+_{\mathrm{even}}([-K,K])$ with respect to the supremum norm.
\end{theorem}

\subsection{Module (A2): Lipschitz continuity}

\begin{theorem}[Lipschitz, Part~I]\label{thm:lipschitz}
For every $K > 0$, there exists $L_Q(K) > 0$ such that
\[
|Q(\Phi_1) - Q(\Phi_2)| \le L_Q(K)\,\|\Phi_1 - \Phi_2\|_\infty
\]
for all admissible test functions $\Phi_1, \Phi_2$ supported in $[-K,K]$.
\end{theorem}

\subsection{Module (A3): Uniform Toeplitz barrier}

\begin{theorem}[Archimedean floor, Part~II]\label{thm:floor}
Fix $t_{\mathrm{sym}} = 3/50$ and $B_{\min} = 3$. Then for every $B \ge B_{\min}$ and every $\theta \in \TT$,
\[
P_A(\theta) \ge c_* := \frac{11}{10},
\]
where $P_A$ is the periodized Archimedean symbol.
\end{theorem}

\subsection{Module (RKHS): Uniform prime cap}

\begin{theorem}[Prime cap, Part~II]\label{thm:prime-cap}
For $t_{\mathrm{rkhs}} \ge 1$,
\[
\|T_P\| \le \rho(t_{\mathrm{rkhs}}) \le \rho(1) < \frac{1}{25} < \frac{c_*}{4},
\]
where $T_P$ is the prime sampling operator.
\end{theorem}

\subsection{Combined: The uniform A3 bridge}

\begin{theorem}[A3 bridge, Part~II]\label{thm:A3-bridge}
For $B \ge B_{\min}$, $t_{\mathrm{sym}} = 3/50$, $t_{\mathrm{rkhs}} \ge 1$, and $M \ge M_0^{\mathrm{unif}}$,
\[
\lambda_{\min}(T_M[P_A] - T_P) \ge \frac{c_*}{4} > 0.
\]
Consequently, $Q(\Phi_{B,t_{\mathrm{sym}}}) \ge 0$ for the associated Fej\'er--heat test function.
\end{theorem}

\subsection{Module summary}

\begin{table}[h]
\centering
\caption{Analytic module stack}
\label{tab:modules}
\begin{tabular}{llll}
\hline
\textbf{Module} & \textbf{Statement} & \textbf{Key constant} & \textbf{Paper} \\
\hline
T0 & Guinand--Weil normalization & --- & III (this paper) \\
A1$'$ & Fej\'er--heat density on $\WW_K$ & --- & I \\
A2 & Lipschitz continuity $L_Q(K)$ & --- & I \\
A3 & Uniform Archimedean floor & $c_* = 11/10$ & II \\
RKHS & Uniform prime cap & $\rho(1) < 1/25$ & II \\
\hline
\end{tabular}
\end{table}

%==============================================================================
\section{Main Closure: Positivity on the Weil Cone}\label{sec:closure}
%==============================================================================

\subsection{The Weil cone}

\begin{definition}[Weil cone]
For $K > 0$, define
\[
\WW_K := \{\Phi \in C^\infty_c(\RR) : \mathrm{supp}\,\Phi \subset [-K,K],\ \Phi(-\xi) = \Phi(\xi),\ \Phi \ge 0\}.
\]
The full Weil cone is $\WW := \bigcup_{K > 0} \WW_K$.
\end{definition}

\subsection{Main positivity theorem}

\begin{theorem}[Main positivity]\label{thm:main-positivity}
Assume modules \textup{(T0)}, \textup{(A1$'$)}, \textup{(A2)}, and the uniform A3 bridge (\Cref{thm:A3-bridge}). Then
\[
Q(\Phi) \ge 0 \qquad \text{for every } \Phi \in \WW.
\]
\end{theorem}

\begin{proof}
Fix $K > 0$. We show $Q(\Phi) \ge 0$ for all $\Phi \in \WW_K$.

\textbf{Step 1: Positivity on generators.} By \Cref{thm:A3-bridge}, $Q(\Phi_{B,t_{\mathrm{sym}}}) \ge 0$ for every Fej\'er--heat generator $\Phi_{B,t_{\mathrm{sym}}}$ with $B \ge B_{\min}$. Since $Q$ is linear and $Q(\Phi) \ge 0$ for generators, we have $Q(g) \ge 0$ for every $g \in \GG_K$.

\textbf{Step 2: Extension by density and continuity.} Let $\Phi \in \WW_K$. By \Cref{thm:density}, there exists a sequence $(g_n) \subset \GG_K$ with $\|g_n - \Phi\|_\infty \to 0$. By \Cref{thm:lipschitz},
\[
|Q(g_n) - Q(\Phi)| \le L_Q(K)\,\|g_n - \Phi\|_\infty \to 0.
\]
Since $Q(g_n) \ge 0$ for all $n$, taking the limit gives $Q(\Phi) \ge 0$.

\textbf{Step 3: Union over compacts.} For any $\Phi \in \WW$, there exists $K > 0$ with $\Phi \in \WW_K$. By Step~2, $Q(\Phi) \ge 0$. Since $\Phi$ was arbitrary, $Q \ge 0$ on $\WW$.
\end{proof}

\begin{remark}[Uniform input only]
The proof uses only the uniform bounds from \Cref{thm:A3-bridge}. No K-dependent schedules or numerical certificates enter the argument.
\end{remark}

%==============================================================================
\section{The Weil Criterion and the Riemann Hypothesis}\label{sec:weil}
%==============================================================================

\subsection{Weil's equivalence}

The following theorem is due to Weil~\cite{Weil1952,Weil1972}; see also~\cite{Guinand1948} for the underlying explicit formula.

\begin{theorem}[Weil's positivity criterion]\label{thm:weil-criterion}
Let $Q$ be the Weil functional attached to the Riemann zeta function $\zeta(s)$ in the normalization of \Cref{sec:T0}, and let $\WW$ be the Weil cone. Then the following are equivalent:
\begin{enumerate}
\item[\textup{(i)}] The Riemann Hypothesis: all nontrivial zeros of $\zeta(s)$ satisfy $\Re s = 1/2$.
\item[\textup{(ii)}] $Q(\Phi) \ge 0$ for every $\Phi \in \WW$.
\end{enumerate}
\end{theorem}

\begin{proof}[Proof sketch]
The Guinand--Weil explicit formula expresses $Q(\Phi)$ as a sum over zeros:
\[
Q(\Phi) = \sum_{\rho} \widehat{\Phi}\Bigl(\frac{\rho - 1/2}{i}\Bigr),
\]
where $\rho$ ranges over the nontrivial zeros. If all zeros lie on $\Re s = 1/2$, then $(\rho - 1/2)/i$ is real, and for even nonnegative $\Phi$ the Fourier transform $\widehat{\Phi}$ evaluated at real points is real; the positivity of $Q$ then follows from properties of the Fourier transform of nonnegative functions.

Conversely, if some zero $\rho_0$ has $\Re \rho_0 \ne 1/2$, one can construct a test function $\Phi$ localized near the imaginary part of $(\rho_0 - 1/2)/i$ that makes $Q(\Phi) < 0$. See~\cite{Weil1972} for the complete argument.
\end{proof}

\subsection{Main theorem}

\begin{theorem}[Riemann Hypothesis]\label{thm:RH}
Assume the analytic module chain:
\begin{enumerate}
\item[\textup{(T0)}] Guinand--Weil normalization (\Cref{prop:GW-match});
\item[\textup{(A1$'$)}] Fej\'er--heat density (\Cref{thm:density});
\item[\textup{(A2)}] Lipschitz continuity (\Cref{thm:lipschitz});
\item[\textup{(A3)}] Uniform Archimedean floor (\Cref{thm:floor});
\item[\textup{(RKHS)}] Uniform prime cap (\Cref{thm:prime-cap}).
\end{enumerate}
Then the Riemann Hypothesis is true: all nontrivial zeros of the Riemann zeta function $\zeta(s)$ lie on the critical line $\Re s = 1/2$.
\end{theorem}

\begin{proof}
By \Cref{thm:main-positivity}, $Q(\Phi) \ge 0$ for every $\Phi \in \WW$. By \Cref{thm:weil-criterion}, this positivity is equivalent to the Riemann Hypothesis.
\end{proof}

\begin{remark}[Scope]
The normalization in (T0) matches the Guinand--Weil conventions, so \Cref{thm:weil-criterion} applies verbatim. No numerical tables or computer-assisted certificates are used in the proof of \Cref{thm:RH}; all bounds are derived from explicit analytic estimates.
\end{remark}

%==============================================================================
\section{Discussion}\label{sec:discussion}
%==============================================================================

\subsection{The modular architecture}

The proof is organized as a chain of independent modules, each verifiable separately:

\begin{center}
\begin{tabular}{ccccc}
\textbf{Paper I} & & \textbf{Paper II} & & \textbf{Paper III} \\
\hline
A1$'$ (density) & $\longrightarrow$ & & & \\
A2 (Lipschitz) & $\longrightarrow$ & & & Main closure \\
& & A3 (Toeplitz) & $\longrightarrow$ & $\Downarrow$ \\
& & RKHS (prime cap) & $\longrightarrow$ & Weil $\Rightarrow$ RH
\end{tabular}
\end{center}

This architecture mirrors standard programmatic proofs in mathematical physics, where theorems are structured around robust intermediate objects and stable limit transitions.

\subsection{What is new}

\begin{enumerate}
\item \textbf{Consistent symbol definition}: The Archimedean symbol $P_A$ is defined once via periodization of the Fej\'er--heat window, and its Lipschitz modulus is proved directly for that same object.

\item \textbf{Uniform prime cap}: The RKHS bound uses a single scale $t_{\mathrm{rkhs}} \ge 1$ independent of $K$, eliminating all K-dependent schedules.

\item \textbf{Two-scale decoupling}: The symbol parameter $t_{\mathrm{sym}}$ and the RKHS parameter $t_{\mathrm{rkhs}}$ are independent; no coupling is required.
\end{enumerate}

\subsection{Relation to other approaches}

The Weil criterion has been studied from various perspectives:

\begin{itemize}
\item \textbf{Li's criterion}~\cite{Li1997} and the \textbf{Jensen polynomial program}~\cite{GriffinOnoRolenZagier2019} provide logically equivalent reformulations of RH.

\item \textbf{Zero-density estimates}~\cite{GuthMaynard2024} encode the zeta problem as a spectral estimate, a viewpoint adopted in our Toeplitz bridge.

\item \textbf{The de Bruijn--Newman constant}~\cite{RodgersTao2020} shows RH may be ``barely true,'' motivating our explicit slack tracking.

\item \textbf{Noncommutative approaches}~\cite{ConnesMarcolli2008} highlight how positivity hinges on operator factorizations.
\end{itemize}

Our approach is distinct: we construct a self-contained, verifiable chain from Toeplitz positivity to Weil positivity, with all critical steps amenable to formal verification.

\subsection{Verification philosophy}

The proof is designed for transparency and auditability:

\begin{enumerate}
\item All constants are explicit and computable.
\item No numerical tables or floating-point computations enter the main chain.
\item Each module is self-contained with clear interfaces.
\item The Lean formalization project (in progress) aims to machine-verify the entire chain.
\end{enumerate}

%==============================================================================
% Bibliography
%==============================================================================

\begin{thebibliography}{99}

\bibitem{ConnesMarcolli2008}
A.~Connes and M.~Marcolli.
\newblock {\em Noncommutative Geometry, Quantum Fields and Motives}.
\newblock Colloquium Publications, vol.~55. American Mathematical Society, 2008.

\bibitem{GriffinOnoRolenZagier2019}
M.~Griffin, K.~Ono, L.~Rolen, and D.~Zagier.
\newblock Jensen polynomials for the {R}iemann zeta function and other sequences.
\newblock {\em Proc.\ Natl.\ Acad.\ Sci.\ USA}, 116(23):11103--11110, 2019.

\bibitem{Guinand1948}
A.~P.~Guinand.
\newblock A summation formula in the theory of prime numbers.
\newblock {\em Proc.\ London Math.\ Soc.\ (2)}, 50:107--119, 1948.

\bibitem{GuthMaynard2024}
L.~Guth and J.~Maynard.
\newblock New large value estimates for {D}irichlet polynomials.
\newblock {\em arXiv:2405.20552}, 2024.

\bibitem{Li1997}
X.-J.~Li.
\newblock The positivity of a sequence of numbers and the {R}iemann hypothesis.
\newblock {\em J.\ Number Theory}, 65(2):325--333, 1997.

\bibitem{Paper1}
E.~Malamutmann.
\newblock Fej\'er--heat generators and {L}ipschitz control for the {W}eil quadratic functional.
\newblock {\em Preprint}, 2025.
\newblock Part~I of this series.

\bibitem{Paper2}
E.~Malamutmann.
\newblock Toeplitz barrier and {RKHS} prime contraction for the {W}eil criterion.
\newblock {\em Preprint}, 2026.
\newblock Part~II of this series.

\bibitem{RodgersTao2020}
B.~Rodgers and T.~Tao.
\newblock The de {B}ruijn--{N}ewman constant is non-negative.
\newblock {\em Forum Math.\ Pi}, 8:e6, 2020.

\bibitem{Weil1952}
A.~Weil.
\newblock Sur les ``formules explicites'' de la th\'eorie des nombres premiers.
\newblock {\em Comm.\ S\'em.\ Math.\ Univ.\ Lund [Medd.\ Lunds Univ.\ Mat.\ Sem.]}, Tome Suppl\'ementaire:252--265, 1952.

\bibitem{Weil1972}
A.~Weil.
\newblock Sur les formules explicites de la th\'eorie des nombres.
\newblock {\em Math.\ USSR Izvestija}, 6(1):1--17, 1972.

\end{thebibliography}

\end{document}
